% GENERATED FILE. Edit canonical lessons, then run npm run generate:books.
% canonical-source-hash: fnv1a64:b2deb2d152984ea5
% canonical-lessons: LA-C487-aquilo, LA-C487-auster, LA-C487-favonius, LA-C487-perflo, LA-C487-serenitas

\chapter{North wind, South wind, West wind, To blow through, Clear weather}
\label{ch:la-north-wind-south-wind-west-wind-to-blow-through-clear-weather}
\hlchaptermodality{487}

\begin{chapteropening}
\textbf{By the end of this chapter:} \emph{I can talk about the winds: the north, south and west winds, the wind blowing across the land, and clear weather after it.}

Complete the last lesson of chapter 487: I can talk about the winds: the north, south and west winds, the wind blowing across the land, and clear weather after it.
\end{chapteropening}

\section[aquilō, aquilōnis]{aquilō, aquilōnis — the north wind; the north}

\label{lesson:LA-C487-aquilo}



\begin{quote}
Before the new one: say the Latin for to be stiff with cold, then the Latin for to feel cold.
\end{quote}

\subsection*{You'll want to know: aquilō, aquilōnis}
\textbf{aquilō, aquilōnis} — \textquotedblleft{}the north wind; the north\textquotedblright{}.

\begin{quote}
Hieme aquilō ventus frīgidus est.
\end{quote}

\begin{quote}
\emph{In winter the north wind is a cold wind.}
\end{quote}

Third declension, masculine: \emph{aquilō, aquilōnis m.} A cold, stormy wind.

\begin{cousinweb}
The origin of \emph{aquilō} is not settled. Lewis \& Short suggest \emph{aqua} 'water', for its rain, or \emph{aquilus} 'dark', for its storm clouds; EtymDB gives no parent.
\end{cousinweb}

\begin{grammarlens}[title={What you've built}]
One more word for telling, explaining and describing at B1.
\end{grammarlens}

\subsection*{Your turn}
\begin{itemize}
\raggedright
  \item \emph{Say it:} \emph{aquilō, aquilōnis}
  \item \emph{Say it:} \emph{aquilō, aquilōnis}, once more
  \item \emph{Say it:} a sentence of your own with \emph{aquilō, aquilōnis}
  \item \emph{Recall:} say the Latin for to be stiff with cold, then the Latin for to feel cold, then say \emph{aquilō, aquilōnis} again
\end{itemize}

\begin{tcolorbox}[breakable,colback=teal!4,colframe=teal!35!black,title={Before you move on}]
What does aquilō, aquilōnis mean? (\textquotedblleft{}The north wind; the north\textquotedblright{}.) Say it once more.
\end{tcolorbox}
\section[auster, austrī]{auster, austrī — the south wind; the south}

\label{lesson:LA-C487-auster}



\begin{quote}
Before the new one: say the Latin for to feel cold, then the Latin for the north wind.
\end{quote}

\subsection*{You'll want to know: auster, austrī}
\textbf{auster, austrī} — \textquotedblleft{}the south wind; the south\textquotedblright{}.

\begin{quote}
Auster calidus nūbēs et pluviam fert.
\end{quote}

\begin{quote}
\emph{The warm south wind brings clouds and rain.}
\end{quote}

Second declension, masculine: \emph{auster, austrī m.} English \emph{austral} comes from it.

\begin{cousinweb}
\emph{auster} is inherited; EtymDB lists English \emph{east} and Latin \emph{aurōra} 'dawn' as cousins, while Lewis \& Short derive it from a root meaning 'to burn', the hot wind. English \emph{austral} comes from it.
\end{cousinweb}

\begin{grammarlens}[title={What you've built}]
One more word for telling, explaining and describing at B1.
\end{grammarlens}

\subsection*{Your turn}
\begin{itemize}
\raggedright
  \item \emph{Say it:} \emph{auster, austrī}
  \item \emph{Say it:} \emph{auster, austrī}, once more
  \item \emph{Say it:} a sentence of your own with \emph{auster, austrī}
  \item \emph{Recall:} say the Latin for to feel cold, then the Latin for the north wind, then say \emph{auster, austrī} again
\end{itemize}

\begin{tcolorbox}[breakable,colback=teal!4,colframe=teal!35!black,title={Before you move on}]
What does auster, austrī mean? (\textquotedblleft{}The south wind; the south\textquotedblright{}.) Say it once more.
\end{tcolorbox}
\section[favōnius, favōniī]{favōnius, favōniī — the west wind (of spring)}

\label{lesson:LA-C487-favonius}



\begin{quote}
Before the new one: say the Latin for the north wind, then the Latin for the south wind.
\end{quote}

\subsection*{You'll want to know: favōnius, favōniī}
\textbf{favōnius, favōniī} — \textquotedblleft{}the west wind (of spring)\textquotedblright{}.

\begin{quote}
Vēre favōnius venit, et flōrēs aperiuntur.
\end{quote}

\begin{quote}
\emph{In spring the west wind comes, and the flowers open.}
\end{quote}

Second declension, masculine: \emph{favōnius, favōniī m.} The mild wind that brings spring.

\begin{cousinweb}
\emph{favōnius} is built on \emph{faveō} 'to be kind, to favour' (EtymDB; Lewis \& Short): the kindly wind of spring. The Alpine wind name \emph{Föhn} is usually traced to it.
\end{cousinweb}

\begin{grammarlens}[title={What you've built}]
One more word for telling, explaining and describing at B1.
\end{grammarlens}

\subsection*{Your turn}
\begin{itemize}
\raggedright
  \item \emph{Say it:} \emph{favōnius, favōniī}
  \item \emph{Say it:} \emph{favōnius, favōniī}, once more
  \item \emph{Say it:} a sentence of your own with \emph{favōnius, favōniī}
  \item \emph{Recall:} say the Latin for the north wind, then the Latin for the south wind, then say \emph{favōnius, favōniī} again
\end{itemize}

\begin{tcolorbox}[breakable,colback=teal!4,colframe=teal!35!black,title={Before you move on}]
What does favōnius, favōniī mean? (\textquotedblleft{}The west wind (of spring)\textquotedblright{}.) Say it once more.
\end{tcolorbox}
\section[perflō, perflāre]{perflō, perflāre — to blow through, to blow across}

\label{lesson:LA-C487-perflo}



\begin{quote}
Before the new one: say the Latin for the south wind, then the Latin for the west wind.
\end{quote}

\subsection*{You'll want to know: perflō, perflāre}
\textbf{perflō, perflāre} — \textquotedblleft{}to blow through, to blow across\textquotedblright{}.

\begin{quote}
Ventus magnus agrōs perflat, et arborēs moventur.
\end{quote}

\begin{quote}
\emph{A strong wind blows across the fields, and the trees move.}
\end{quote}

\emph{perflō, perflāre, perflāvī, perflātum}, with an accusative of the place. The simple verb \emph{flō} 'to blow' gives English \emph{inflate} through \emph{īnflāre}.

\begin{cousinweb}
\emph{perflō} is built from \emph{per} 'through' and \emph{flō} 'to blow' (Lewis \& Short). \emph{flō} is inherited, and EtymDB lists English \emph{blow} as a cousin; English \emph{inflate} and \emph{flatulent} come from \emph{flō}.
\end{cousinweb}

\begin{grammarlens}[title={What you've built}]
One more word for telling, explaining and describing at B1.
\end{grammarlens}

\subsection*{Your turn}
\begin{itemize}
\raggedright
  \item \emph{Say it:} \emph{perflō, perflāre}
  \item \emph{Say it:} \emph{perflō, perflāre}, once more
  \item \emph{Say it:} a sentence of your own with \emph{perflō, perflāre}
  \item \emph{Recall:} say the Latin for the south wind, then the Latin for the west wind, then say \emph{perflō, perflāre} again
\end{itemize}

\begin{tcolorbox}[breakable,colback=teal!4,colframe=teal!35!black,title={Before you move on}]
What does perflō, perflāre mean? (\textquotedblleft{}To blow through, to blow across\textquotedblright{}.) Say it once more.
\end{tcolorbox}
\section[serēnitās, serēnitātis]{serēnitās, serēnitātis — clear weather, fair weather; calm}

\label{lesson:LA-C487-serenitas}



\begin{quote}
Before the new one: say the Latin for the west wind, then the Latin for to blow through.
\end{quote}

\subsection*{You'll want to know: serēnitās, serēnitātis}
\textbf{serēnitās, serēnitātis} — \textquotedblleft{}clear weather, fair weather; calm\textquotedblright{}.

\begin{quote}
Post ventum serēnitās est, et caelum caeruleum est.
\end{quote}

\begin{quote}
\emph{After the wind the weather is clear, and the sky is blue.}
\end{quote}

Third declension, feminine: \emph{serēnitās, serēnitātis f.}, from \emph{serēnus} 'clear'. English \emph{serenity} comes from it.

\begin{cousinweb}
\emph{serēnitās} is built on \emph{serēnus} 'clear, cloudless' with the suffix \emph{-itās} (EtymDB; Lewis \& Short). English \emph{serenity} comes from it.
\end{cousinweb}

\begin{grammarlens}[title={What you've built}]
One more word for telling, explaining and describing at B1.
\end{grammarlens}

\subsection*{Your turn}
\begin{itemize}
\raggedright
  \item \emph{Say it:} \emph{serēnitās, serēnitātis}
  \item \emph{Say it:} \emph{serēnitās, serēnitātis}, once more
  \item \emph{Say it:} a sentence of your own with \emph{serēnitās, serēnitātis}
  \item \emph{Recall:} say the Latin for the west wind, then the Latin for to blow through, then say \emph{serēnitās, serēnitātis} again
\end{itemize}

\begin{tcolorbox}[breakable,colback=teal!4,colframe=teal!35!black,title={Before you move on}]
What does serēnitās, serēnitātis mean? (\textquotedblleft{}Clear weather, fair weather; calm\textquotedblright{}.) Say it once more.
\end{tcolorbox}
